\chapter{Conclusão}
\label{chap:conclusao}

Este trabalho propôs uma arquitetura unificada de dupla via para a classificação de imagens histopatológicas sob restrição amostral, integrando a dinâmica topológica multiescala à capacidade representacional de modelos profundos de visão computacional. As imagens foram modeladas como redes complexas ponderadas e direcionadas em raios de conexão crescentes, gerando sequências de atributos de conectividade tecidual que foram codificadas por modelos neurais recorrentes. Paralelamente, modelos de visão computacional pré-treinados extraíram atributos densos diretamente dos pixels. As representações latentes de ambas as modalidades foram equalizadas dimensionalmente por projeção linear e normalização por camadas, alimentando um classificador conjunto ajustado de ponta a ponta. Essa formulação estabeleceu uma alternativa aos descritores tradicionais de escala fixa ao viabilizar a exploração de relações teciduais em diferentes raios de conexão, indicando benefícios preditivos decorrentes da integração multimodal em determinadas tarefas de classificação.

Os experimentos conduzidos sobre o conjunto de teste independente responderam à questão de pesquisa e delimitaram o alcance das hipóteses formuladas a partir de três observações centrais. Primeiro, a via topológica apresentou capacidade discriminativa própria, inclusive quando utilizada isoladamente, como atestado pelo desempenho dos modelos sequenciais no conjunto multiclasse NHL e pela separação quase perfeita nas bases hepáticas LA e LG. Segundo, a integração das modalidades demonstrou benefícios mais evidentes na tarefa multiclasse NHL, caracterizada por elevada sobreposição citológica entre subtipos de linfoma, na qual a arquitetura de dupla via superou a via visual isolada em 1,6 ponto percentual no F1 macro ($0{,}938 \pm 0{,}026$) e reduziu a dispersão entre sementes em 30\%. Terceiro, a contribuição da fusão mostrou dependência direta das características da tarefa, não sendo necessariamente vantajosa em todos os conjuntos avaliados: em tarefas binárias com marcante contraste morfológico (CR e LG), as representações visuais puras já atingiram o patamar máximo de desempenho, tornando a fusão neutra, ao passo que em conjuntos com reduzidíssimo volume de teste, a exemplo do UCSB, a maior parametrização da rede integrada introduziu leve sensibilidade a variações estocásticas.

O objetivo geral e os objetivos específicos foram integralmente cumpridos ao longo das etapas metodológicas. As contribuições deste estudo compreendem o estabelecimento de uma abordagem que trata medidas de redes complexas como séries ordenadas por escala acopladas à visão profunda, a comprovação da eficiência paramétrica da via topológica, que acrescentou apenas 189 mil parâmetros (1,0\% da rede) sem consumo de memória gráfica, e a caracterização comparativa detalhada de seis codificadores sequenciais e seis modelos visuais em cinco bases histopatológicas, delimitando o impacto da fusão frente a cada modalidade isolada.

Para trabalhos futuros, as investigações foram delineadas em dois eixos principais: o aprimoramento da arquitetura e a ampliação da validação e interpretabilidade. No eixo voltado à arquitetura, a concatenação simples de vetores pode ser substituída por mecanismos de atenção cruzada para ponderar de forma adaptativa as escalas topológicas e os atributos visuais, somada à investigação de resoluções flexíveis de entrada para mitigar o descarte de bordas teciduais decorrente do recorte central. No eixo referente à validação e interpretabilidade, as propostas contemplam a extensão da abordagem para lâminas inteiras digitalizadas por meio de aprendizado fracamente supervisionado por múltiplas instâncias, a validação em biópsias de novos órgãos sob escassez de dados e o desenvolvimento de mapas de explicabilidade associados às medidas nodais das redes complexas.
